\chapter{अवकलन: पहला पाठ्यक्रम}\label{ch:g11:deriv}

कोई \omterm{def:g11:func:function}{फलन} कितनी तेज़ी से बदलता है? किसी एक बिंदु पर इसका उत्तर \emph{\omterm{def:g11:deriv:derivative}{अवकलज}} है,
जो पूरे व्यावहारिक गणित की सबसे उपयोगी अकेली संख्या है। यह अध्याय उसे \omterm{def:g11:deriv:rate}{औसत
परिवर्तन दरों} से खड़ा करता है, \omterm{def:g11:func:reference}{आधार फलनों} के लिए उसे निकालता है, और उसके
चिह्न से किसी \omterm{def:g11:func:function}{फलन} के परिवर्तन पढ़ता है। यहाँ की हर बात \cref{ch:g12:deriv} में और आगे बढ़ाई
तथा गहरी की जाती है।

\section{औसत परिवर्तन दर}

\begin{definition}[औसत परिवर्तन दर]\label{def:g11:deriv:rate}
मान लीजिए $f$ किसी \omterm{def:g10:numbers:interval}{अंतराल} $I$ पर परिभाषित है और $I$ में $a \neq b$ हैं। $a$ और
$b$ के बीच $f$ की \emph{औसत परिवर्तन दर}\index{औसत परिवर्तन दर} है
\[
\frac{f(b) - f(a)}{b - a}.
\]
यह वक्र के बिंदुओं $\bigl(a, f(a)\bigr)$ और $\bigl(b, f(b)\bigr)$ से होकर जाने वाली \emph{छेदक
रेखा}\index{छेदक रेखा} की \omterm{def:g10:reffunc:affine}{प्रवणता} है।
\end{definition}

\begin{example}
एक कार दोपहर 2 बजे से 4 बजे के बीच $140$ km चलती है: उसकी औसत चाल $\frac{140}{2} = 70$ km/h
है। पर गति-मापक हर क्षण एक \emph{तात्क्षणिक} चाल दिखाता है --- और \omterm{def:g11:deriv:derivative}{अवकलज} ठीक
इसी धारणा को सुनिश्चित कर देता है।
\end{example}

\section{किसी बिंदु पर अवकलज}

\begin{definition}[किसी बिंदु पर अवकलज]\label{def:g11:deriv:derivative}
मान लीजिए $f$ ऐसे \omterm{def:g10:numbers:interval}{अंतराल} पर परिभाषित है जिसमें $a$ है। $h \neq 0$ के लिए $a$
और $a + h$ के बीच \omterm{def:g11:deriv:rate}{औसत परिवर्तन दर} बनाइए:
\[
r(h) = \frac{f(a+h) - f(a)}{h}.
\]
यदि $h$ के $0$ के जितना चाहें उतना पास आने पर $r(h)$ किसी एक नियत संख्या के
पास पहुँचता जाए, तो हम कहते हैं कि $f$
\emph{$a$ पर अवकलनीय}\index{अवकलज}\index{अवकलनीय} है; यह संख्या \emph{$a$
पर $f$ का अवकलज} है, जिसे $f'(a)$ लिखते हैं:
\[
f'(a) = \lim_{h \to 0} \frac{f(a+h) - f(a)}{h}.
\]
\end{definition}

\begin{example}\label{ex:g11:deriv:atpoint}
मान लीजिए $f(x) = x^2$ और $a = 1$ हैं। $h \neq 0$ के लिए:
\[
\frac{f(1+h) - f(1)}{h}
= \frac{(1+h)^2 - 1}{h}
= \frac{1 + 2h + h^2 - 1}{h}
= \frac{h(2 + h)}{h}
= 2 + h .
\]
जैसे-जैसे $h$ $0$ के पास आता है, यह $2$ के पास पहुँचता है: \omterm{def:g11:func:function}{फलन} $x^2$ $1$
पर \omterm{def:g11:deriv:derivative}{अवकलनीय} है और $f'(1) = 2$। रणनीति पर ध्यान दीजिए: \emph{भागफल को तब तक सरल
कीजिए जब तक $h$ किसी हर में न रह जाए}, और तभी $h \to 0$ कीजिए।
\end{example}

\begin{example}[एक अनवकलनीय बिंदु]
मान लीजिए $f(x) = \abs{x}$ और $a = 0$ हैं। तब $\frac{f(h) - f(0)}{h} = \frac{\abs h}{h}$, जो $h > 0$ के लिए $1$ और $h < 0$ के लिए
$-1$ के बराबर है: कोई एक संख्या नहीं आती, और $\abs x$ $0$ पर \omterm{def:g11:deriv:derivative}{अवकलनीय} नहीं है।
वहाँ उसके वक्र में एक \emph{कोना} है।
\end{example}

\section{स्पर्श रेखा}

\begin{definition}[स्पर्श रेखा]\label{def:g11:deriv:tangent}
यदि $f$ $a$ पर \omterm{def:g11:deriv:derivative}{अवकलनीय} हो, तो बिंदु $A = \bigl(a, f(a)\bigr)$ पर वक्र की \emph{स्पर्श
रेखा}\index{स्पर्श रेखा} वह रेखा है जो $A$ से होकर जाती है और जिसकी \omterm{def:g10:reffunc:affine}{प्रवणता}
$f'(a)$ है। उसका \omterm{def:g10:algebra:equation}{समीकरण} है
\[
y = f(a) + f'(a)(x - a).
\]
\end{definition}

\omterm{def:g11:deriv:tangent}{स्पर्श रेखा} $A$ से होकर जाने वाली छेदक रेखाओं की सीमांत स्थिति है: जैसे-जैसे
$h \to 0$, दूसरा बिंदु $\bigl(a+h, f(a+h)\bigr)$ वक्र पर सरककर $A$ की ओर आता है, और छेदक की \omterm{def:g10:reffunc:affine}{प्रवणता}
$\frac{f(a+h)-f(a)}{h}$ $f'(a)$ के पास पहुँचती है।

\begin{omfigure}
\begin{tikzpicture}
\begin{axis}[omaxis,
  width=8.8cm, height=6cm,
  xmin=-0.4, xmax=3.2, ymin=-0.6, ymax=4.3,
  xtick={1,2.2}, xticklabels={$a$,$a+h$}, ytick=\empty]
  \addplot[omProp, thick, domain=-0.1:2.9] {1 + 1.6*(x-1)};
  \addplot[omThm, thick, domain=-0.35:2.6] {1 + 1*(x-1)};
  \addplot[omDef, thick, domain=-0.4:2.75] {x^2/2 + 0.5};
  \fill (axis cs:1,1) circle (1.5pt) node[above left, font=\small] {$A$};
  \fill (axis cs:2.2,2.92) circle (1.5pt)
    node[above left, font=\small] {$B$};
  \node[omProp, font=\small, anchor=west] at (axis cs:2.62,3.75)
    {छेदक रेखा};
  \node[omThm, font=\small, anchor=west] at (axis cs:2.62,2.45)
    {स्पर्श रेखा};
\end{axis}
\end{tikzpicture}

{\small $A$ और $B = \bigl(a+h, f(a+h)\bigr)$ से होकर जाने वाली \omterm{def:g11:deriv:rate}{छेदक रेखा} की \omterm{def:g10:reffunc:affine}{प्रवणता} $\frac{f(a+h)-f(a)}{h}$ है।
जैसे-जैसे $B$ सरककर $A$ की ओर आता है, छेदक झुककर \omterm{def:g11:deriv:tangent}{स्पर्श रेखा} बन जाती है,
जिसकी \omterm{def:g10:reffunc:affine}{प्रवणता} $f'(a)$ है।}
\end{omfigure}

\begin{example}\label{ex:g11:deriv:tangenteq}
$a = 1$ पर $f(x) = x^2$ के लिए: $f(1) = 1$ और $f'(1) = 2$ (\cref{ex:g11:deriv:atpoint}), इसलिए \omterm{def:g11:deriv:tangent}{स्पर्श रेखा} है
\[
y = 1 + 2(x - 1) = 2x - 1 .
\]
यही वह रेखा $d_2$ है जो \cref{exo:g11:quad:11} में शुद्ध बीजगणित से मिली थी।
\end{example}

\section{आधार फलनों के अवकलज}

\begin{proposition}\label{prop:g11:deriv:refderivatives}
अपने-अपने अवकलनीयता के \omterm{def:g11:func:function}{प्रांत} पर:
\[
(c)' = 0, \quad
(x)' = 1, \quad
(x^2)' = 2x, \quad
(x^3)' = 3x^2, \quad
\left(\frac1x\right)' = -\frac{1}{x^2}, \quad
(\sqrt x\,)' = \frac{1}{2\sqrt x} \ (x > 0).
\]
और अधिक सामान्य रूप से, हर \omterm{def:g10:numbers:sets}{पूर्णांक} $n \geq 1$ के लिए $(x^n)' = n\,x^{n-1}$।
\end{proposition}

\begin{proof}
हर उपपत्ति \cref{ex:g11:deriv:atpoint} की रणनीति का अनुसरण करती है: दर को सरल कीजिए, फिर $h \to 0$
कीजिए। \omterm{def:g11:func:function}{प्रांत} में $a$ स्थिर रखिए।

\emph{अचर और तत्समक।} सभी $h$ के लिए $\frac{c - c}{h} = 0$ और $\frac{(a+h) - a}{h} = 1$: सीमाएँ $0$ और $1$
हैं।

\emph{वर्ग।} $\frac{(a+h)^2 - a^2}{h} = \frac{2ah + h^2}{h} = 2a + h
\to 2a$।

\emph{घन।} $(a+h)^3 = a^3 + 3a^2h + 3ah^2 + h^3$, इसलिए दर $3a^2 + 3ah + h^2 \to 3a^2$ है।

\emph{व्युत्क्रम।} $a \neq 0$ के लिए और इतने छोटे $h$ के लिए कि $a + h \neq 0$ हो:
\[
\frac{1}{h}\left(\frac{1}{a+h} - \frac{1}{a}\right)
= \frac{1}{h} \cdot \frac{a - (a+h)}{a(a+h)}
= \frac{-1}{a(a+h)}
\xrightarrow[h \to 0]{} -\frac{1}{a^2}.
\]

\emph{वर्गमूल।} $a > 0$ के लिए संयुग्मी से गुणा कीजिए:
\[
\frac{\sqrt{a+h} - \sqrt a}{h}
= \frac{(a + h) - a}{h\,(\sqrt{a+h} + \sqrt a)}
= \frac{1}{\sqrt{a+h} + \sqrt a}
\xrightarrow[h \to 0]{} \frac{1}{2\sqrt a}.
\]

\emph{सामान्य घात।} ढर्रा $2a$, $3a^2$ आगे भी चलता है: $(a+h)^n$ का प्रसार करने
पर $a^n + n\,a^{n-1}h$ के बाद के सभी पद $h^2$ उठाए रहते हैं, इसलिए दर $n\,a^{n-1}$ की ओर जाती है।
(पूरा आगमन \cref{ch:g12:deriv} में दिया गया है।)
\end{proof}

\section{अवकलन के नियम}

\begin{theorem}[संक्रियाएँ]\label{thm:g11:deriv:rules}
मान लीजिए $u$ और $v$ $I$ पर \omterm{def:g11:deriv:derivative}{अवकलनीय} हैं और $\lambda \in \R$ है। तब $u + v$, $\lambda u$,
$uv$ $I$ पर \omterm{def:g11:deriv:derivative}{अवकलनीय} हैं, और जहाँ-जहाँ $v \neq 0$ है वहाँ $\frac{u}{v}$ भी, तथा
\[
(u+v)' = u' + v', \qquad
(\lambda u)' = \lambda u', \qquad
(uv)' = u'v + uv', \qquad
\left(\frac{u}{v}\right)' = \frac{u'v - uv'}{v^2}.
\]
\end{theorem}

\begin{proof}[पहले तीन नियमों की उपपत्ति]
$a \in I$ स्थिर रखिए और संक्षेप में $\Delta u = u(a+h) - u(a)$, $\Delta v = v(a+h) - v(a)$ लिखिए।

\emph{योग।} $u + v$ की दर $\frac{\Delta u + \Delta v}{h} = \frac{\Delta u}{h} + \frac{\Delta v}{h}
\to u'(a) + v'(a)$ है।

\emph{गुणज।} $\lambda u$ की दर $\lambda \frac{\Delta u}{h} \to \lambda u'(a)$ है।

\emph{गुणनफल।} मिश्रित पद $u(a+h)\,v(a)$ जोड़िए और घटाइए:
\begin{align*}
\frac{u(a+h)v(a+h) - u(a)v(a)}{h}
&= \frac{u(a+h)v(a+h) - u(a+h)v(a) + u(a+h)v(a) - u(a)v(a)}{h}\\
&= u(a+h)\,\frac{\Delta v}{h} + v(a)\,\frac{\Delta u}{h}.
\end{align*}
जैसे-जैसे $h \to 0$: $\frac{\Delta v}{h} \to v'(a)$, $\frac{\Delta u}{h} \to
u'(a)$, और $u(a+h) \to u(a)$ (उसकी दर की सीमा है, इसलिए $\Delta u = h \cdot \frac{\Delta u}{h} \to 0$)।
अतः दर $u(a)v'(a) + v(a)u'(a)$ की ओर जाती है।

भागफल का नियम इन पंक्तियों पर सिद्ध होता है; उसे \emph{इस स्तर पर स्वीकार
कर लिया जाता है} और \cref{ch:g12:deriv} में सिद्ध किया जाता है।
\end{proof}

\begin{example}\label{ex:g11:deriv:computations}
\emph{बहुपद।} $f(x) = 3x^3 - 5x^2 + 7x - 2$:
\[
f'(x) = 3 \times 3x^2 - 5 \times 2x + 7 = 9x^2 - 10x + 7 .
\]

\emph{भागफल।} $g(x) = \frac{2x+1}{x^2+1}$: $u = 2x + 1$, $v = x^2
+ 1$ के साथ
\[
g'(x) = \frac{2(x^2+1) - (2x+1)\,2x}{(x^2+1)^2}
= \frac{-2x^2 - 2x + 2}{(x^2+1)^2}.
\]

\emph{गुणनफल।} $x > 0$ के लिए $k(x) = x\sqrt x$: $k'(x) = 1 \times \sqrt x + x \times \frac{1}{2\sqrt x}
= \sqrt x + \frac{\sqrt x}{2} = \frac{3\sqrt x}{2}$।
\end{example}

\section{अवकलज का चिह्न और परिवर्तन}

\begin{theorem}[अवकलज और परिवर्तन]\label{thm:g11:deriv:variations}
मान लीजिए $f$ किसी \omterm{def:g10:numbers:interval}{अंतराल} $I$ पर \omterm{def:g11:deriv:derivative}{अवकलनीय} है।
\begin{enumerate}
  \item यदि $I$ पर $f' \geq 0$ हो, तो $f$ $I$ पर \omterm{def:g11:func:monotone}{वर्धमान} है।
  \item यदि $I$ पर $f' \leq 0$ हो, तो $f$ $I$ पर \omterm{def:g10:functions:variations}{ह्रासमान} है।
  \item यदि $I$ पर $f' > 0$ हो (सिवाय अधिक से अधिक सीमित बिंदुओं के जहाँ वह
        लुप्त होता है), तो $f$ $I$ पर निरंतर \omterm{def:g11:func:monotone}{वर्धमान} है।
\end{enumerate}
\end{theorem}

\begin{proof}
अंतर्ज्ञान स्पष्ट है --- जिस वक्र की सभी \omterm{def:g11:deriv:tangent}{स्पर्श रेखाएँ} ऊपर की ओर इशारा करती
हों उसे चढ़ना ही होगा --- पर कठोर उपपत्ति के लिए माध्यमान प्रमेय चाहिए। यह
परिणाम \emph{इस स्तर पर स्वीकार कर लिया जाता है}; \cref{thm:g12:deriv:variations} देखिए।
\end{proof}

\begin{proposition}[चरम मान]\label{prop:g11:deriv:extrema}
यदि $f'$ $a$ पर $+$ से $-$ में चिह्न बदले, तो $f$ का $a$ पर एक
स्थानीय \omterm{def:g10:functions:extrema}{अधिकतम} है; और $-$ से $+$ में बदले तो स्थानीय \omterm{def:g10:functions:extrema}{न्यूनतम}। ऐसे बिंदु
पर $f'(a) = 0$ होता है और \omterm{def:g11:deriv:tangent}{स्पर्श रेखा} क्षैतिज होती है।
\end{proposition}

\begin{proof}
यदि $a$ पर समाप्त होने वाले किसी \omterm{def:g10:numbers:interval}{अंतराल} पर $f' \geq 0$ हो और $a$ से आरंभ होने
वाले किसी \omterm{def:g10:numbers:interval}{अंतराल} पर $f' \leq 0$ हो, तो $f$ $f(a)$ तक बढ़ता है और उसके बाद घटता है
(\cref{thm:g11:deriv:variations}): आसपास $f(a)$ सबसे बड़ा है। चूँकि $f'$ $a$ पर चिह्न बदलता है और
वहाँ परिभाषित है, इसलिए $f'(a) = 0$।
\end{proof}

\begin{method}[किसी फलन का अध्ययन]\label{met:g11:deriv:study}
\begin{enumerate}
  \item $f'(x)$ निकालिए और उसे सरल कीजिए --- सबसे अच्छा हो कि गुणनखंडित रूप
        में;
  \item हर \omterm{def:g10:numbers:interval}{अंतराल} पर $f'(x)$ का चिह्न तय कीजिए (द्विघात $f'$ के लिए \cref{thm:g11:quad:sign}
        का उपयोग कीजिए);
  \item परिणामों को \emph{\omterm{def:g10:functions:table}{परिवर्तन-सारणी}} में दर्ज कीजिए: $f'$ के चिह्न की
        एक पंक्ति, $f$ के तीरों की एक पंक्ति, और मोड़ के बिंदुओं पर $f$
        के मान;
  \item निष्कर्ष निकालिए: चरम मान, $f(x) = k$ के हलों की संख्या, असमिकाएँ।
\end{enumerate}
\end{method}

\begin{example}\label{ex:g11:deriv:study}
$\R$ पर $f(x) = x^3 - 3x + 1$ का अध्ययन कीजिए। पहले $f'(x) = 3x^2 - 3 = 3(x-1)(x+1)$: यह एक द्विघात है जिसके \omterm{def:g11:quad:discriminant}{मूल} $-1$
और $1$ हैं, और जो उनके बाहर धनात्मक है। इसलिए $f$ $\intoc{-\infty}{-1}$ पर \omterm{def:g11:func:monotone}{वर्धमान},
$\intcc{-1}{1}$ पर \omterm{def:g10:functions:variations}{ह्रासमान} और $\intco{1}{+\infty}$ पर \omterm{def:g11:func:monotone}{वर्धमान} है, तथा उसका स्थानीय \omterm{def:g10:functions:extrema}{अधिकतम} $f(-1) = -1 + 3 + 1 = 3$ और
स्थानीय \omterm{def:g10:functions:extrema}{न्यूनतम} $f(1) = 1 - 3 + 1 = -1$ है।
\end{example}

\begin{omfigure}
\begin{tikzpicture}
\begin{axis}[omaxis,
  width=8.4cm, height=5.8cm,
  xmin=-2.6, xmax=2.6, ymin=-3.6, ymax=4.4,
  xtick={-1,1}, ytick={-1,3}]
  \addplot[omProp, thick, domain=-1.7:-0.3] {3};
  \addplot[omProp, thick, domain=0.3:1.7] {-1};
  \addplot[omDef, thick, domain=-2.25:2.25, samples=90] {x^3 - 3*x + 1};
  \fill (axis cs:-1,3) circle (1.5pt);
  \fill (axis cs:1,-1) circle (1.5pt);
\end{axis}
\end{tikzpicture}

{\small $f(x) = x^3 - 3x + 1$ का वक्र (\cref{ex:g11:deriv:study}): स्थानीय \omterm{def:g10:functions:extrema}{अधिकतम} $(-1, 3)$ और स्थानीय \omterm{def:g10:functions:extrema}{न्यूनतम} $(1, -1)$
पर क्षैतिज \omterm{def:g11:deriv:tangent}{स्पर्श रेखाएँ} (नारंगी), जहाँ $f'$ चिह्न बदलता है।}
\end{omfigure}

\begin{example}[इष्टतमीकरण]\label{ex:g11:deriv:box}
$12 \times 12$ cm की एक चादर के हर कोने से भुजा $x$ का वर्ग काटकर और मोड़कर खुला
डिब्बा बनाया जाता है। $0 < x < 6$ के लिए उसका आयतन $V(x) = x(12 - 2x)^2$ है। प्रसार करने पर $V(x) = 4x^3 - 48x^2 + 144x$,
इसलिए
\[
V'(x) = 12x^2 - 96x + 144 = 12(x^2 - 8x + 12) = 12(x - 2)(x - 6).
\]
$\intoo{0}{6}$ पर गुणनखंड $x - 6$ ऋणात्मक है, इसलिए $V'(x)$ का चिह्न $-(x-2)$ जैसा है:
$2$ से पहले धनात्मक, बाद में ऋणात्मक। आयतन $x = 2$ cm पर \omterm{def:g10:functions:extrema}{अधिकतम} होता है, और
वह $V(2) = 2 \times 8^2 = 128$ cm$^3$ है।
\end{example}

\section{अभ्यास}

\begin{exercise}[$\star$]\label{exo:g11:deriv:1}
परिभाषा का उपयोग करके (परिवर्तन दर, फिर $h \to 0$) $f(x) = x^2$ के लिए $f'(2)$ निकालिए,
फिर $g(x) = x^2 + 3x$ के लिए $g'(1)$।
\end{exercise}

\begin{exercise}[$\star$]\label{exo:g11:deriv:2}
अवकलन कीजिए:
\[
f(x) = 4x^3 - 2x^2 + x - 7, \qquad
g(x) = 5\sqrt{x} + \frac{3}{x}, \qquad
h(x) = (2x+1)(x^2 - 3) .
\]
\end{exercise}

\begin{exercise}[$\star$]\label{exo:g11:deriv:3}
\omterm{def:g10:coordgeom:system}{भुज} $a$ वाले बिंदु पर $f$ के वक्र की \omterm{def:g11:deriv:tangent}{स्पर्श रेखा} का \omterm{def:g10:algebra:equation}{समीकरण} बताइए:
\[
f(x) = x^2 - 3x, \ a = 2; \qquad
f(x) = \frac1x, \ a = 1; \qquad
f(x) = \sqrt x, \ a = 4 .
\]
\end{exercise}

\begin{exercise}[$\star$]\label{exo:g11:deriv:4}
भागफलों का अवकलन कीजिए:
\[
f(x) = \frac{x+2}{x-1}, \qquad
g(x) = \frac{x^2}{x^2+1}, \qquad
h(x) = \frac{1}{x^2 + x + 1} .
\]
\end{exercise}

\begin{exercise}[$\star\star$]\label{exo:g11:deriv:5}
$\R$ पर $f(x) = x^3 - 6x^2 + 9x - 2$ के परिवर्तन का अध्ययन कीजिए (\omterm{def:g10:functions:table}{परिवर्तन-सारणी}, स्थानीय चरम
मान)।
\end{exercise}

\begin{exercise}[$\star\star$]\label{exo:g11:deriv:6}
$\intoo{-1}{+\infty}$ पर $f(x) = \frac{x^2 + 3}{x + 1}$ के परिवर्तन का अध्ययन कीजिए।
\end{exercise}

\begin{exercise}[$\star\star$]\label{exo:g11:deriv:7}
मान लीजिए $f(x) = x^2 - 4x + 5$ है।
\begin{enumerate}
  \item वक्र का वह बिंदु ज्ञात कीजिए जहाँ \omterm{def:g11:deriv:tangent}{स्पर्श रेखा} क्षैतिज है।
  \item वे बिंदु ज्ञात कीजिए जहाँ \omterm{def:g11:deriv:tangent}{स्पर्श रेखा} रेखा $y = 2x + 1$ के समांतर है।
\end{enumerate}
\end{exercise}

\begin{exercise}[$\star\star$]\label{exo:g11:deriv:8}
एक दीवार के सहारे $60$ m बाड़ से (तीन भुजाएँ) आयताकार बाड़ा बनाया जाता है।
क्षेत्रफल को चौड़ाई $x$ के \omterm{def:g11:func:function}{फलन} के रूप में व्यक्त कीजिए, और \omterm{def:g11:deriv:derivative}{अवकलज} से
\omterm{def:g10:functions:extrema}{अधिकतम} क्षेत्रफल वाली विमाएँ ज्ञात कीजिए। (\cref{ch:g11:quad} की शीर्ष वाली विधि से तुलना
कीजिए।)
\end{exercise}

\begin{exercise}[$\star\star$]\label{exo:g11:deriv:9}
$\intoo{0}{+\infty}$ पर \omterm{def:g11:func:function}{फलन} $f(x) = \frac{x+1}{2} - \sqrt x$ का अध्ययन करके दिखाइए कि सभी $x > 0$ के लिए $\sqrt{x} \leq \frac{x + 1}{2}$ है।
\end{exercise}

\begin{exercise}[$\star\star$]\label{exo:g11:deriv:10}
एक बेलनाकार डिब्बे में $500$ cm$^3$ समाना चाहिए। उसका पृष्ठीय क्षेत्रफल (दो
वृत्त और पार्श्व) $S = 2\pi r^2 + 2\pi r h$ है, जहाँ $\pi r^2 h = 500$ है। $S$ को केवल $r$ के \omterm{def:g11:func:function}{फलन} के
रूप में व्यक्त कीजिए और दिखाइए कि $S'(r) = 4\pi r -
\frac{1000}{r^2}$। इससे वह त्रिज्या निकालिए जो
पृष्ठीय क्षेत्रफल को \omterm{def:g10:functions:extrema}{न्यूनतम} करती है, यथार्थ व्यंजक के रूप में।
\end{exercise}

\begin{exercise}[$\star\star\star$]\label{exo:g11:deriv:11}
बाहरी बिंदु $P(1, -3)$ से होकर \omterm{def:g11:quad:parabola}{परवलय} $y = x^2$ की कितनी \omterm{def:g11:deriv:tangent}{स्पर्श रेखाएँ} जाती हैं? उनके
\omterm{def:g10:algebra:equation}{समीकरण} ज्ञात कीजिए। ($a$ को स्पर्श बिंदु का \omterm{def:g10:coordgeom:system}{भुज} मानिए और यह व्यक्त कीजिए
कि $a$ पर की \omterm{def:g11:deriv:tangent}{स्पर्श रेखा} $P$ से होकर जाती है।)
\end{exercise}

\begin{exercise}[$\star\star\star$]\label{exo:g11:deriv:12}
मान लीजिए $f(x) = x^3 - 3x + 1$ है (\cref{ex:g11:deriv:study} देखिए)। \omterm{def:g10:functions:table}{परिवर्तन-सारणी} का उपयोग करके \omterm{def:g10:numbers:sets}{वास्तविक
संख्या} $k$ के मान के अनुसार \omterm{def:g10:algebra:equation}{समीकरण} $f(x) = k$ के हलों की संख्या ज्ञात कीजिए।
\end{exercise}

\section{समस्या: स्पर्श रेखा की तीन नौकरियाँ}

\begin{problem}[{सप्ताहांत समस्या --- परवलयिक दर्पण की उपपत्ति, न्यूटन की
मूल-खोजी मशीन, और लट्ठे में छिपा सबसे मज़बूत शहतीर}]
\label{pb:g11:deriv:1}
\omterm{def:g11:deriv:derivative}{अवकलज} \omterm{def:g11:deriv:tangent}{स्पर्श रेखाएँ} खींचता है; सुनने में यह मामूली प्रतिभा लगती है। यह समस्या
दिखाती है कि \omterm{def:g11:deriv:tangent}{स्पर्श रेखा} एक साथ तीन नौकरियाँ करती है: वह \cref{pb:g11:quad:1} में लटकी छोड़
दी गई हेडलाइट वाली प्रमेय सिद्ध करती है (\omterm{def:g11:quad:parabola}{परवलय} के अक्ष के समांतर आने वाली हर
किरण \omterm{pb:g11:quad:1}{नाभि} से होकर परावर्तित होती है), वह आम उपयोग के सबसे तेज़ समीकरण-हल
कलनविधि को चलाती है --- जो हर कैलकुलेटर और हर तंत्रिका जाल के भीतर है --- और
वह ज्ञात करती है कि किसी गोल लट्ठे से सबसे मज़बूत आयताकार शहतीर कौन-सा निकल
सकता है।

\medskip
\noindent\textbf{भाग I --- \omterm{def:g11:deriv:tangent}{स्पर्श रेखा} में प्रवाह।}
\begin{enumerate}
  \item $f(x) = 3x^2 - 5x + 1$ और $g(x) = x^3 - 12x$ का अवकलन कीजिए (\cref{thm:g11:deriv:rules}); और $x = 1$ पर $\left(x^2
        \right)'$ को
        परिभाषा से (परिवर्तन दर, फिर $h \to 0$) फिर से निकालिए।
  \item $a = 2$ पर $y = x^3 - 12x$ की \omterm{def:g11:deriv:tangent}{स्पर्श रेखा} ज्ञात कीजिए। उसमें ख़ास क्या है, और
        वह किस बात का संकेत है (\cref{prop:g11:deriv:extrema})?
  \item बिंदु $(a, a^2)$ पर \omterm{def:g11:quad:parabola}{परवलय} $y = x^2$ की सामान्य \omterm{def:g11:deriv:tangent}{स्पर्श रेखा} स्थापित कीजिए:
        \[
        y = 2ax - a^2 .
        \]
  \item वह \omterm{def:g11:deriv:tangent}{स्पर्श रेखा} $y$-अक्ष को कहाँ काटती है? इससे रेखाचित्रकार का
        नुस्ख़ा निकालिए: \omterm{def:g11:quad:parabola}{परवलय} पर ऊँचाई $h$ वाले बिंदु $P$ पर \omterm{def:g11:deriv:tangent}{स्पर्श
        रेखा} खींचने के लिए $P$ को मूल बिंदु से इतनी गहराई\,\dots\ वाले
        अक्ष-बिंदु से जोड़िए। उसे लिखिए।
  \item $g(x) = x^3 - 12x$ की पूरी \omterm{def:g10:functions:table}{परिवर्तन-सारणी} बनाइए ($g'$ का चिह्न, स्थानीय चरम मान
        और उनके मान, \cref{met:g11:deriv:study})।
\end{enumerate}

\medskip
\noindent\textbf{भाग II --- दर्पण प्रमेय।}
मान लीजिए $P(a, a^2)$ \omterm{def:g11:quad:parabola}{परवलय} $y = x^2$ का एक बिंदु है ($a \neq 0$ लीजिए), $F\left(0, \frac14\right)$ \omterm{pb:g11:quad:1}{नाभि} है, और
$T$ वह बिंदु है जहाँ $P$ पर की \omterm{def:g11:deriv:tangent}{स्पर्श रेखा} $y$-अक्ष को काटती है
(प्रश्न 4)। \cref{pb:g11:quad:1} से याद कीजिए कि $PF = a^2 + \frac14$।
\begin{enumerate}[resume]
  \item दूरी $FT$ निकालिए।
  \item निष्कर्ष निकालिए कि त्रिभुज $FPT$ $F$ पर समद्विबाहु है। अतः
        कौन-से दो कोण बराबर हैं?
  \item $P$ से होकर अक्ष के समांतर आने वाली किरण रेखा $(TF)$ के समांतर है
        (दोनों ऊर्ध्वाधर)। एकांतर कोणों का उपयोग करके दिखाइए कि $P$ पर की
        \omterm{def:g11:deriv:tangent}{स्पर्श रेखा} आती हुई ऊर्ध्वाधर किरण और रेखाखंड $[PF]$ के साथ
        \emph{बराबर कोण} बनाती है।
  \item भौतिकी: प्रकाश किसी वक्र से वैसे ही परावर्तित होता है जैसे उसकी
        \omterm{def:g11:deriv:tangent}{स्पर्श रेखा} से, और उसी कोण पर लौटता है जिस पर आया था (परावर्तन का
        नियम)। \emph{दर्पण प्रमेय} का निष्कर्ष निकालिए: अक्ष के समांतर आने
        वाली हर किरण \omterm{pb:g11:quad:1}{नाभि} से होकर परावर्तित होती है --- और समय उलट दें तो
        \omterm{pb:g11:quad:1}{नाभि} पर रखा बल्ब समांतर प्रकाश फेंकता है। \cref{pb:g11:quad:1} की तश्तरी और
        हेडलाइट अब प्रमेय हैं।
  \item $a = 1$ पर समद्विबाहु त्रिभुज की संख्यात्मक जाँच कीजिए: $T$, $FP$
        और $FT$ निकालिए।
\end{enumerate}

\medskip
\noindent\textbf{भाग III --- न्यूटन की मशीन।}
$f(x) = 0$ हल करने के लिए न्यूटन ने यह सुझाया: किसी अनुमान $x_0$ से शुरू कीजिए,
$x_0$ पर की \omterm{def:g11:deriv:tangent}{स्पर्श रेखा} के साथ नीचे अक्ष तक जाइए, और उसका $x$-अंतःखंड अगला
अनुमान मान लीजिए।
\begin{enumerate}[resume]
  \item विधि का सूत्र निकालिए:
        \[
        x_{n + 1} = x_n - \frac{f(x_n)}{f'(x_n)} .
        \]
  \item इसे $x_0 = 1$ से $f(x) = x^2 - 2$ पर चलाइए: $x_1$, $x_2$, $x_3$ को यथार्थ भिन्नों
        के रूप में निकालिए। आपने अभी दो हज़ार साल पुराना कौन-सा अनुक्रम ---
        और \cref{pb:g10:functions:1} की कौन-सी मशीन --- फिर से खड़ा कर दिया है? इस संयोग के
        पीछे की सर्वसमिका सिद्ध कीजिए:
        \[
        x - \frac{x^2 - 2}{2x} = \frac12\left(x +
        \frac2x\right).
        \]
  \item $x^3 = 100$ हल कीजिए (\cref{pb:g10:reffunc:1} का बचत-योजनाओं वाला पार-बिंदु): $x_0 = 5$ से
        $f(x) = x^3 - 100$ पर न्यूटन के दो क़दम लगाइए (चार दशमलव स्थान), और कैलकुलेटर के
        $\sqrt[3]{100}$ से तुलना कीजिए।
  \item तोड़फोड़: विधि को $x_0 = 2$ पर $f(x) = x^3 - 12x$ से चलाने का प्रयास कीजिए। क्या
        गड़बड़ होती है, और क्यों (प्रश्न 2)? व्यावहारिक चेतावनी बताइए।
  \item प्रश्न 12 में सही दशमलव स्थान मोटे तौर पर $1 \to 3 \to 6$ की तरह बढ़े। इसकी
        तुलना द्विभाजन से कीजिए (हर क़दम पर \omterm{def:g10:numbers:interval}{अंतराल} आधा करना) और एक वाक्य में
        बताइए कि कैलकुलेटर, GPS अभिग्राही और तंत्रिका जालों का प्रशिक्षण
        सभी न्यूटन के विचार पर क्यों चलते हैं।
\end{enumerate}

\medskip
\noindent\textbf{भाग IV --- सबसे मज़बूत शहतीर।}
चौड़ाई $w$ और गहराई $d$ वाला एक आयताकार शहतीर व्यास $30$ cm के गोल
लट्ठे से चीरा जाता है, इसलिए $w^2 + d^2 = 900$। शहतीर की दृढ़ता $w d^2$ के समानुपाती है
(गहराई दो बार गिनी जाती है: कड़ियाँ किनारे पर खड़ी रखी जाती हैं!)।
\begin{enumerate}[resume]
  \item दृढ़ता को केवल $w$ के \omterm{def:g11:func:function}{फलन} के रूप में व्यक्त कीजिए: $S(w) = 900w - w^3$, और किस
        \omterm{def:g10:numbers:interval}{अंतराल} पर?
  \item अवकलन कीजिए, \omterm{def:g10:functions:table}{परिवर्तन-सारणी} बनाइए, और इष्टतम $w$ तथा $d$ ज्ञात
        कीजिए (यथार्थ मान, फिर मिलीमीटर तक)।
  \item दिखाइए कि इष्टतम अनुपात $d = w\sqrt2$ को संतुष्ट करते हैं --- वही पुराना
        बढ़इयों का अनुपात। (अतिरिक्त जानकारी: लट्ठे के व्यास को तीन बराबर
        भागों में बाँटकर वृत्त तक लंब खड़े करने से ठीक यही आयत बन जाता है।)
  \item भोला चुनाव \emph{वर्गाकार} शहतीर है ($w = d$)। उसकी और इष्टतम शहतीर
        की दृढ़ता निकालिए, और बढ़ई का लाभ प्रतिशत में बताइए।
  \item समापन --- \omterm{def:g11:deriv:tangent}{स्पर्श रेखा} की तीन नौकरियाँ, एक-एक वाक्य में: ज्यामिति
        (दर्पण प्रमेय), संख्यात्मक गणना (न्यूटन की मशीन), इष्टतमीकरण
        (शहतीर) --- और तीनों को एक ही वस्तु चलाती है। और आगे का संकेत:
        कक्षा 12 उत्तलता जोड़ती है --- वक्र के सापेक्ष \omterm{def:g11:deriv:tangent}{स्पर्श रेखा} किस ओर
        है --- और तीनों को और पैना कर देती है।
\end{enumerate}
\end{problem}
